\par\noindent\textbf{(1)} 令 \(X=(x_1+x_2)/2\) 及 \(x=x_1-x_2\)，则 \(T=m\dot X^2+\frac14m\dot x^2\)。相对运动的约化质量为 \(m_r=m/2\)，Hamilton量 \(H=P_X^2/(4m)+p_x^2/(2m_r)+kx^2/2+\beta x^4/4\)，质心动量守恒。
\par\noindent\textbf{(2)} Euler方程 \(m_r\ddot x+kx+\beta x^3=0\)。定义线性角频率 \(\omega_0^2=k/m_r=2k/m\)，则 \(\ddot x+\omega_0^2x+(\beta/m_r)x^3=0\)。正四次非线性使振幅越大回复力越强，属于硬化型。
\par\noindent\textbf{(3)} 设 \(x\approx A\cos(\Omega t)\)，用 \(\cos^3u=(3\cos u+\cos3u)/4\)，丢弃三次谐波的首阶作用得 \(\Omega^2\approx\omega_0^2+3\beta A^2/(4m_r)\)。因此 \(\Omega\approx\omega_0[1+3\beta A^2/(8k)]\)，要求 \(\beta A^2/k\ll1\)。
\par\noindent\textbf{(4)} 相对运动能量 \(E_r=\frac12m_r\dot x^2+\frac12kx^2+\frac14\beta x^4\) 守恒。因 \(k,\beta>0\)，势能随 \(|x|\to\infty\) 发散，给定有限能量的相对轨道有界；若换 \(\beta<0\)，四次截断势在大振幅不再有界，不能据小振幅展开预测全局稳定。
